% !TeX root = BookN.tex

\title{Boundary Element Method for 2D Fracture Mechanics Analysis in Quasicrystalline Materials}
\titlerunning{Boundary Element Method for 2D Fracture Mechanics Analysis}

\author{
	Roman Kushnir,
	Iaroslav Pasternak,
	Heorhiy Sulym, and
	Ihor Hotsyk
}

\institute{
	Roman Kushnir 
	\at
	\AffPidstryhach
    \\
	\email{dyrector@iapmm.lviv.ua}
	\and
	Iaroslav Pasternak 
	\at
	\AffVolyn
	\\
	\email{iaroslav.pasternak@vnu.edu.ua}
	\and
	Heorhiy Sulym 
	\at
	\AffPidstryhach
	\\
	\email{gtsulym@gmail.com}
	\and
	Ihor Hotsyk 
	\at
	\AffVolyn
	\\
	\email{gotsyk.igor@vnu.edu.ua}
}

\maketitle

\graphicspath{{24_Kushnir/}}

\abstract{ An efficient boundary element formulation is developed for the analysis
	of quasicrystalline bodies containing cracks under thermomechanical
	loading. A specialized numerical technique is proposed for the accurate
	computation of phonon--phason stress intensity factors near crack tips.
	The approach demonstrates high accuracy and numerical stability, and
	the obtained results are validated through comparison with available
	analytical solutions. It enables a detailed investigation of the effect
	of crack size on generalized stress intensity factors in finite
	domains. The developed framework also provides new insight into
	thermomechanical phenomena specific to quasicrystals.
}

%\keywords{quasicrystals, boundary element method, phonon--phason stress intensity factors, thermomechanical loading, fracture mechanics}

\section{Introduction}

Quasicrystals constitute a unique class of solid materials characterized
by a non-periodic yet long-range ordered atomic structure. Since their
discovery in the early 1980s \cite{Kushnir1}, quasicrystals have attracted
considerable attention owing to their unusual physical properties and
their potential applications in various technological fields. One of
the most notable features of quasicrystalline materials is their
exceptionally low thermal conductivity, which makes them suitable for
high-temperature applications and energy-efficient technologies
\cite{Kushnir40}.

Unlike conventional crystals, quasicrystals exhibit symmetries that
are forbidden in periodic lattices, such as fivefold and icosahedral
symmetries. This structural peculiarity leads to distinctive physical
behavior, including low thermal conductivity \cite{Kushnir40}, high
hardness \cite{Kushnir41}, and significant resistance to wear and
corrosion \cite{Kushnir42}. The aperiodic atomic arrangement impedes
phonon propagation and thus reduces the ability of the material to
conduct heat. As a result, quasicrystals are promising candidates for
thermal barrier coatings and insulating layers in aerospace and power
generation technologies.

The brittle behavior of quasicrystals is commonly attributed to the
absence of conventional dislocation-mediated plasticity
\cite{Kushnir40}. Instead, deformation is often governed by defects
specific to quasiperiodic structures, such as phasons. However,
phason-related processes are typically slow and inefficient in
relaxing stress concentrations near crack tips. Consequently, cracks
in quasicrystals may propagate abruptly once a critical stress level
is reached \cite{Kushnir39}. Therefore, given their numerous potential
applications, studying the fracture mechanics of quasicrystals under
combined thermal and mechanical loading is of significant importance.

In recent years, substantial progress has been made in developing
analytical and numerical approaches for fracture analysis in cracked
quasicrystals. Fan et al.~\cite{Kushnir5} investigated crack behavior
in one-dimensional hexagonal quasicrystals by reducing the governing
problem to the determination of four harmonic potentials. Pi et
al.~\cite{Kushnir7} applied complex variable techniques to plane
electroelastic problems in 6mm one-dimensional quasicrystal solids.
Fan and Mai~\cite{Kushnir8} introduced a Lekhnitskii-type method for
plane elasticity problems in quasicrystals, whereas Gao et
al.~\cite{Kushnir9} generalized the Stroh formalism for
one-dimensional quasicrystal media. These approaches were further
extended to icosahedral and decagonal quasicrystals by Li and Liu
\cite{Kushnir10,Kushnir11}. A Stroh-like formalism was also employed
by Radi and Mariano~\cite{Kushnir12} for crack analysis, while Zhang
et al.~\cite{Kushnir13} derived fundamental solutions for
one-dimensional piezoelectric quasicrystals within the same framework.

Comparatively fewer studies have addressed plane thermoelasticity
problems in quasicrystals. In \cite{Kushnir6}, an axisymmetric
thermoelastic problem was analyzed to determine thermal stresses in
a homogeneous disk composed of a cubic quasicrystal. Guo et
al.~\cite{Kushnir22} applied a Stroh-type formalism to solve the plane
thermoelastic problem for a two-dimensional quasicrystal containing
an elliptic hole. The stress singularity induced by thermal effects
in one-dimensional hexagonal piezoelectric quasicrystal composites
was examined by Mu et al.~\cite{Kushnir23} using an extended Stroh
formalism. Lu et al.~\cite{Kushnir38} employed Kachanov's method to
investigate crack interaction in one-dimensional thermoelastic
quasicrystals. Pasternak et al.~\cite{Kushnir34} developed a general
boundary integral equation approach for thermoelastic problems in
cracked quasicrystals.

However, to the best of the authors' knowledge, a general boundary
element formulation for fracture mechanics analysis in thermoelastic
quasicrystals has not yet been developed. To address this gap, the
present study builds upon the boundary integral equations proposed
by Pasternak et al.~\cite{Kushnir34} and develops a boundary element
method for two-dimensional fracture analysis in thermoelastic
quasicrystals with general material configurations, including
one-, two-, and three-dimensional quasicrystalline structures.

\section{Governing Relations and Boundary Integral Equations}

Since thermal effects may significantly influence the fracture
behavior of quasicrystals, they must be properly taken into account
in fracture mechanics analyses. In this study, steady-state plane heat
conduction and static plane thermoelasticity problems are considered
for quasicrystalline materials with a general three-dimensional
internal structure.

Following the formulations in \cite{Kushnir6,Kushnir21}, quasicrystal
solids are assumed to obey the classical Fourier law for steady-state
heat conduction,
\begin{equation}
	h_i = -k_{ij}\theta_{,j},
	\label{KushnirEq1}
\end{equation}
where $\theta$ denotes the temperature change with respect to a
reference temperature and $k_{ij}$ are the thermal conductivity
coefficients. The conductivity tensor is assumed to be symmetric,
i.e., $k_{ij}=k_{ji}$.

The heat flux vector field ${\bf h}$ must satisfy the stationary heat
balance equation
\begin{equation}
	\mathrm{div}\,{\bf h} - f_h \equiv h_{i,i} - f_h = 0,
	\label{KushnirEq2}
\end{equation}
where $f_h$ denotes the volumetric heat source density.

Following the formulation presented in \cite{Kushnir34}, the
thermoelastic constitutive relations for quasicrystalline materials
can be written in a unified generalized form as
\begin{equation}
	\tilde{\sigma}_{Ij}
	=
	\tilde{C}_{IjKm}\tilde{u}_{K,m}
	-
	\tilde{\beta}_{Ij}\theta,
	\label{KushnirEq3}
\end{equation}
together with the equilibrium equations
\begin{equation}
	\tilde{\sigma}_{Ij,j}=-\tilde{f}_{I}.
	\label{KushnirEq4}
\end{equation}

In this formulation, all field variables are expressed in terms of
generalized quantities (denoted by tildes), which simultaneously
describe both phonon and phason contributions. These quantities are
defined as
\begin{equation}
	\begin{array}{l}
		\tilde{\sigma}_{ij}=\sigma_{ij}, \qquad
		\tilde{\sigma}_{(i+3)j}=H_{ij}, \qquad
		\tilde{u}_{i}=u_{i}, \qquad
		\tilde{u}_{i+3}=w_{i}, \\[3pt]
		\tilde{C}_{ijkm}=C_{ijkm}, \qquad
		\tilde{C}_{ij(k+3)m}=R_{ijkm}, \\[3pt]
		\tilde{C}_{(i+3)jkm}=R_{kmij}, \qquad
		\tilde{C}_{(i+3)j(k+3)m}=K_{ijkm} \qquad
		(i,j,k,m=1,2,3).
	\end{array}
	\label{KushnirEq5}
\end{equation}
%
Here $\sigma_{ij}$ is the phonon stress tensor (symmetric),
$H_{ij}$ is the phason stress tensor,
$\varepsilon_{ij}=\tfrac12(u_{i,j}+u_{j,i})$ is the phonon strain tensor
(symmetric), and $w_{i,j}$ is the phason strain tensor.
Furthermore, $C_{ijkm}$ denotes the tensor of phonon elastic constants,
$K_{ijkm}$ is the tensor of phason elastic constants, and
$R_{ijkm}$ together with $R'_{ijkm}=R_{kmij}$ describe the
phonon–phason coupling. Finally, $f_i$ and $g_i$ denote the components
of the phonon and phason body force vectors, respectively.

The generalized thermal moduli of quasicrystals are defined as
\begin{equation}
	\begin{array}{rcl}
		\tilde{\beta}_{ij}&=&C_{ijkm}\alpha_{km}+R_{ijkm}\alpha'_{km},\\[4pt]
		\tilde{\beta}_{(i+3)j}&=&R_{kmij}\alpha_{km}+K_{ijkm}\alpha'_{km},
	\end{array}
	\label{KushnirEq6}
\end{equation}
where $\alpha_{ij}$ and $\alpha'_{ij}$ denote the phonon and phason
components of the linear thermal expansion tensor, respectively.
Based on symmetry considerations of the quasicrystalline structure,
the phason thermal expansion components are commonly assumed to vanish,
i.e., $\alpha'_{ij}\equiv0$ (see \cite{Kushnir20} for further details).
In the present work, only weak coupling between the thermal and
mechanical fields is considered.

Throughout this work, the Einstein summation convention is adopted,
where repeated indices imply summation over the range $1$ to $3$.
Partial differentiation with respect to spatial coordinates is
indicated by a subscript comma. For example,
\[
u_{i,j}=\frac{\partial u_i}{\partial x_j}.
\]

Using complex variable methods in combination with the extended Stroh
formalism, boundary integral equations for the Stroh complex potentials
were derived in \cite{Kushnir34}. By employing the relationships
between these complex potentials and the physical fields -- such
as temperature, heat flux, extended displacement, and extended stress
-- the corresponding boundary integral equations for plane heat conduction
and plane thermoelasticity in quasicrystals can be formulated.

Thus, the boundary integral equations for plane heat conduction of
a cracked quasicrystal solid are as follows.
\begin{itemize}
	\item For a collocation point $\mathbf{y}\in\Gamma$, assuming the boundary
	$\Gamma$ of a solid is smooth at $\mathbf{y}$, the following boundary
	integral equation holds:
\end{itemize}
\begin{multline}
	\frac{1}{2}\theta(\mathbf{y})
	=
	\mathrm{RPV}\!\int_{\Gamma}
	\Theta^{*}(\mathbf{x},\mathbf{y})\,h_n(\mathbf{x})\,d\Gamma(\mathbf{x})
	-
	\mathrm{CPV}\!\int_{\Gamma}
	H^{*}(\mathbf{x},\mathbf{y})\,\theta(\mathbf{x})\,d\Gamma(\mathbf{x})
	\\
	+
	\int_{\Gamma_C}
	\left[
	\Theta^{*}(\mathbf{x},\mathbf{y})\,\Sigma h_n(\mathbf{x})
	-
	H^{*}(\mathbf{x},\mathbf{y})\,\Delta\theta(\mathbf{x})
	\right] d\Gamma(\mathbf{x})
	\\
	-
	\iint_{S}
	\Theta^{*}(\mathbf{x},\mathbf{y})\,f_h(\mathbf{x})\,dS(\mathbf{x}).
	\label{KushnirEq7}
\end{multline}
\begin{itemize}
	
	\item For a collocation point $\mathbf{y}_{C}\in\Gamma_{C}$ located on
	a smooth crack surface $\Gamma_{C}$, the boundary integral equation
	takes the form
	\begin{multline}
		\frac{1}{2}\Delta h_n(\mathbf{y}_C)
		=
		n_i^{+}(\mathbf{y}_C)
		\bigg[
		\mathrm{CPV}\!\int_{\Gamma_C}
		\Theta_i^{**}(\mathbf{x},\mathbf{y}_C)
		\,\Sigma h_n(\mathbf{x})\,d\Gamma(\mathbf{x})
		\\
		-
		\mathrm{HPV}\!\int_{\Gamma_C}
		H_i^{**}(\mathbf{x},\mathbf{y}_C)
		\,\Delta\theta(\mathbf{x})\,d\Gamma(\mathbf{x})
		\\
		+
		\int_{\Gamma}
		\big(
		\Theta_i^{**}(\mathbf{x},\mathbf{y}_C)h_n(\mathbf{x})
		-
		H_i^{**}(\mathbf{x},\mathbf{y}_C)\theta(\mathbf{x})
		\big)
		\,d\Gamma(\mathbf{x})
		\\
		-
		\iint_{S}
		\Theta_i^{**}(\mathbf{x},\mathbf{y}_C)
		f_h(\mathbf{x})\,dS(\mathbf{x})
		\bigg].
		\label{KushnirEq8}
	\end{multline}
	
\end{itemize}

Boundary integral equations for plane thermoelasticity of quasicrystals
can be written as follows.

\begin{itemize}
	
	\item For a collocation point $\mathbf{y}\in\Gamma$, assuming that the
	boundary $\Gamma$ is smooth at $\mathbf{y}$, the boundary integral
	representation of the extended displacement field is
\begin{multline}
	\frac{1}{2}\tilde{u}_I(\mathbf{y})
	=
	\mathrm{RPV}\!\int_{\Gamma}
	U_{IJ}(\mathbf{x},\mathbf{y})\tilde{t}_J(\mathbf{x})\,d\Gamma(\mathbf{x})
	-
	\mathrm{CPV}\!\int_{\Gamma}
	T_{IJ}(\mathbf{x},\mathbf{y})\tilde{u}_J(\mathbf{x})\,d\Gamma(\mathbf{x})
	\\
	+
	\mathrm{RPV}\!\int_{\Gamma}
	R_I(\mathbf{x},\mathbf{y})\theta(\mathbf{x})\,d\Gamma(\mathbf{x})
	+
	\int_{\Gamma}
	V_I(\mathbf{x},\mathbf{y})h_n(\mathbf{x})\,d\Gamma(\mathbf{x})
	\\
	+
	\int_{\Gamma_C}
	\big(
	U_{IJ}(\mathbf{x},\mathbf{y})\Sigma\tilde{t}_J(\mathbf{x})
	-
	T_{IJ}(\mathbf{x},\mathbf{y})\Delta\tilde{u}_J(\mathbf{x})
	\big)\,d\Gamma(\mathbf{x})
	\\
	+
	\int_{\Gamma_C}
	\big(
	R_I(\mathbf{x},\mathbf{y})\Delta\theta(\mathbf{x})
	+
	V_I(\mathbf{x},\mathbf{y})\Sigma h_n(\mathbf{x})
	\big)\,d\Gamma(\mathbf{x})
	\\
	+
	\iint_{S}
	U_{IJ}(\mathbf{x},\mathbf{y})\tilde{f}_J(\mathbf{x})\,dS(\mathbf{x})
	-
	\iint_{S}
	V_I(\mathbf{x},\mathbf{y})f_h(\mathbf{x})\,dS(\mathbf{x}).
	\label{KushnirEq9}
\end{multline}
	
	\item For a collocation point $\mathbf{y}_C\in\Gamma_C$, assuming that
	the crack surface $\Gamma_C$ is smooth at $\mathbf{y}_C$, the boundary
	integral equation for the extended stress discontinuity is
	\begin{multline}
		\frac{1}{2}\Delta\tilde{t}_I(\mathbf{y}_C)
		=
		n_j^{+}(\mathbf{y}_C)
		\bigg[
		\mathrm{CPV}\!\int_{\Gamma_C}
		D_{IjK}(\mathbf{x},\mathbf{y}_C)
		\,\Sigma\tilde{t}_K(\mathbf{x})\,d\Gamma(\mathbf{x})
		\\
		-
		\mathrm{HPV}\!\int_{\Gamma_C}
		S_{IjK}(\mathbf{x},\mathbf{y}_C)
		\,\Delta\tilde{u}_K(\mathbf{x})\,d\Gamma(\mathbf{x})
		+
		\mathrm{CPV}\!\int_{\Gamma_C}
		Q_{Ij}(\mathbf{x},\mathbf{y}_C)
		\,\Delta\theta(\mathbf{x})\,d\Gamma(\mathbf{x})
		\\
		+
		\mathrm{RPV}\!\int_{\Gamma_C}
		W_{Ij}(\mathbf{x},\mathbf{y}_C)
		\,\Sigma h_n(\mathbf{x})\,d\Gamma(\mathbf{x})
		\\
		+
		\int_{\Gamma}
		\big(
		D_{IjK}(\mathbf{x},\mathbf{y}_C)\tilde{t}_K(\mathbf{x})
		-
		S_{IjK}(\mathbf{x},\mathbf{y}_C)\tilde{u}_K(\mathbf{x})
		\big)
		\,d\Gamma(\mathbf{x})
		\\
		+
		\int_{\Gamma}
		\big(
		Q_{Ij}(\mathbf{x},\mathbf{y}_C)\theta(\mathbf{x})
		+
		W_{Ij}(\mathbf{x},\mathbf{y}_C)h_n(\mathbf{x})
		\big)
		\,d\Gamma(\mathbf{x})
		\\
		+
		\iint_{S}
		D_{IjK}(\mathbf{x},\mathbf{y}_C)
		\,\tilde{f}_K(\mathbf{x})\,dS(\mathbf{x})
		-
		\iint_{S}
		W_{Ij}(\mathbf{x},\mathbf{y}_C)
		\,f_h(\mathbf{x})\,dS(\mathbf{x})
		\bigg].
		\label{KushnirEq10}
	\end{multline}
	%
\end{itemize}

	In (\ref{KushnirEq9}–\ref{KushnirEq12}) and throughout the text, RPV denotes the Riemann
	principal value of an improper integral, CPV denotes the Cauchy
	principal value of a weakly singular integral, and HPV denotes the
	Hadamard principal value (or finite part) of a hypersingular integral. The quantity
	$\tilde{t}_{I}=\tilde{\sigma}_{Ij}n_{j}$
	represents the extended surface traction vector, where $n_{j}$ is
	the outward unit normal vector to the boundary $\Gamma$ or to the
	crack surface $\Gamma_{C}$. The superscript ``+'' indicates evaluation
	on the positive face of the crack.
	
	The symbols $\Sigma$ and $\Delta$ denote, respectively, the sum
	\[
	\Sigma f = f^{+}+f^{-}
	\]
	and the jump (difference)
	\[
	\Delta f = f^{+}-f^{-}
	\]
	of a field quantity across the crack surface $\Gamma_{C}$, where
	$f^{+}$ and $f^{-}$ represent the values on the upper and lower crack
	faces.
	

	
	The kernels appearing in the boundary integral equations are defined as follows.
	
	\begin{equation}
		\begin{array}{rclrcl}
			\Theta^{*}(\mathbf{x},\boldsymbol{\xi}) 
			&=& \dfrac{1}{2\pi k_t}
			\ln\left|Z_t(\mathbf{x}-\boldsymbol{\xi})\right|,
			\quad &
			H^{*}(\mathbf{x},\boldsymbol{\xi})
			&=& -k_{ij}\Theta^{*}_{,j}(\mathbf{x},\boldsymbol{\xi})\,n_i(\mathbf{x}),
			\\[1em]
			\Theta_i^{**}(\mathbf{x},\boldsymbol{\xi})
			&=& -k_{ij}\dfrac{\partial}{\partial \xi_j}
			\Theta^{*}(\mathbf{x},\boldsymbol{\xi}),
			\quad &
			H_i^{**}(\mathbf{x},\boldsymbol{\xi})
			&=& -k_{ij}\dfrac{\partial}{\partial \xi_j}
			H^{*}(\mathbf{x},\boldsymbol{\xi}).
		\end{array}
		\label{KushnirEq11}
	\end{equation}
	
	\begin{align}
		\notag
		U_{IJ}(\mathbf{x},\boldsymbol{\xi})
		&=
		\frac{1}{\pi}
		\Im\!\left[
		A_{I\alpha}A_{J\alpha}
		\ln Z_\alpha(\mathbf{x}-\boldsymbol{\xi})
		\right],
		\\[.2em]
		\notag
		T_{IQ}(\mathbf{x},\boldsymbol{\xi})
		&=
		\tilde{C}_{IjKm}\tilde{U}_{KQ,m}(\mathbf{x},\boldsymbol{\xi})\,n_j
		\\[.2em]
		\notag
		k_{pq}V_{I,pq}(\mathbf{x},\boldsymbol{\xi})
		&=
		\tilde{\beta}_{Jk}U_{IJ,k}(\mathbf{x},\boldsymbol{\xi}),
		\\[.2em]
		\notag
		R_I(\mathbf{x},\boldsymbol{\xi})
		&=
		k_{pq}V_{I,q}(\mathbf{x},\boldsymbol{\xi})\,n_p(\mathbf{x})
		\\
		\label{KushnirEq12}
		D_{IjK}(\mathbf{x},\boldsymbol{\xi})
		&=
		\tilde{C}_{IjMp}
		\frac{\partial}{\partial\xi_p}
		U_{MK}(\mathbf{x},\boldsymbol{\xi}),
		\\
		\notag
		S_{IjK}(\mathbf{x},\boldsymbol{\xi})
		&=
		\tilde{C}_{IjMp}
		\frac{\partial}{\partial\xi_p}
		T_{MK}(\mathbf{x},\boldsymbol{\xi})
			\\
			\notag
		Q_{Ij}(\mathbf{x},\boldsymbol{\xi})
		&=
		\tilde{C}_{IjMp}
		\frac{\partial}{\partial\xi_p}
		R_M(\mathbf{x},\boldsymbol{\xi})
		+
		\tilde{\beta}_{Ij}
		H^{*}(\mathbf{x},\boldsymbol{\xi})
		\\
		\notag
		W_{Ij}(\mathbf{x},\boldsymbol{\xi})
		&=
		\tilde{C}_{IjMp}
		\frac{\partial}{\partial\xi_p}
		V_M(\mathbf{x},\boldsymbol{\xi})
		-
		\tilde{\beta}_{Ij}
		\Theta^{*}(\mathbf{x},\boldsymbol{\xi}).
	\end{align}
	%
	Here,
	\[
	k_t=\sqrt{k_{11}k_{22}-k_{12}^{2}},
	\]
	and
	\[
	Z_t(\mathbf{x}) = x_1 + p_t x_2
	\]
	is a complex variable introduced for the heat conduction problem.
	The parameter $p_t$ is the root with a positive imaginary part of
	the characteristic equation for anisotropic heat conduction,
	\[
	k_{22}p_t^2 + 2k_{12}p_t + k_{11} = 0 .
	\]
	
	In the thermoelastic case, $A_{I\alpha}$ and $p_{\alpha}$ are,
	respectively, the eigenvectors and eigenvalues of the extended Stroh
	eigenvalue problem \cite{Kushnir34}. The quantity
	\[
	Z_{\alpha}(\mathbf{x}) = x_{1} + p_{\alpha}x_{2}
	\]
	denotes the complex argument associated with the $\alpha$th eigenvalue.
	
	Boundary conditions provide only half of the boundary data required
	for the field functions $\theta$, $h_{n}$, $\tilde{u}_{I}$, and
	$\tilde{t}_{I}$. Therefore, the boundary integral equations for heat
	conduction~\eqref{KushnirEq7} and  \eqref{KushnirEq8} are solved to determine
	either the unknown temperature $\theta$ (when the heat flux $h_{n}$
	is prescribed on the boundary $\Gamma$) or the unknown normal heat
	flux $h_{n}$ (when the boundary temperature $\theta$ is specified).
	Similarly, on the crack surface $\Gamma_{C}$ these equations are used
	to determine either the temperature discontinuity $\Delta\theta$
	(for a thermally insulated crack) or the heat flux discontinuity
	$\Sigma h_{n}$ (for a thermally conductive crack).
	
	The boundary integral equations for thermoelasticity
    \eqref{KushnirEq9} and \eqref{KushnirEq10} are solved in an analogous manner.
	Specifically, they are used to determine the unknown extended
	displacement $\tilde{u}_{I}$ (when the surface tractions $\tilde{t}_{I}$
	are prescribed) or the unknown extended traction $\tilde{t}_{I}$ (when
	the displacements $\tilde{u}_{I}$ are prescribed) on the boundary
	$\Gamma$. On the crack surface $\Gamma_{C}$, the equations determine
	the displacement discontinuity $\Delta\tilde{u}_{I}$. In the present
	study, symmetric crack loading is assumed, meaning that the traction
	is continuous across the crack faces, i.e.,
	\[
	\Sigma\tilde{t}_{I}\equiv 0 .
	\]

\section{Field Singularity at the Tips of Incisions}

In \cite{Kushnir34}, it was shown that in the most general case of
three-dimensional quasicrystals both the heat flux and the extended
stress fields exhibit a square-root singularity near the crack tips.
Such behavior is characteristic of fracture problems and is consistent
with classical results of linear elastic fracture mechanics, extended
to complex materials such as quasicrystals.

The explicit expressions for the extended stress components near the
crack tip are given by \cite{Kushnir34}
\begin{equation}
	\begin{array}{rcl}
		\tilde{\boldsymbol{\sigma}}_{1}=[\tilde{\sigma}_{I1}]
		&=&
		-\dfrac{1}{\sqrt{2\pi}}
		\Re\!\left\{
		\mathbf{B}
		\left\langle p_{*}Z_{*}^{-1/2}\right\rangle
		\mathbf{B}^{-1}
		\tilde{\mathbf{k}}^{(1)}
		\right\}
		+ O(1),
		\\[1em]
		\tilde{\boldsymbol{\sigma}}_{2}=[\tilde{\sigma}_{I2}]
		&=&
		\dfrac{1}{\sqrt{2\pi}}
		\Re\!\left\{
		\mathbf{B}
		\left\langle Z_{*}^{-1/2}\right\rangle
		\mathbf{B}^{-1}
		\tilde{\mathbf{k}}^{(1)}
		\right\}
		+ O(1),
	\end{array}
	\label{KushnirEq13}
\end{equation}
%
where
\[
\begin{array}{rcl}
\left\langle \sqrt{Z_{*}} \right\rangle
&=&
\mathrm{diag}\!\left[
\sqrt{Z_{1}},\sqrt{Z_{2}},\ldots,\sqrt{Z_{6}}
\right],
\\[1em]
\left\langle \sqrt{p_{*}Z_{*}} \right\rangle
&=&
\mathrm{diag}\!\left[
\sqrt{p_{1}Z_{1}},\sqrt{p_{2}Z_{2}},\ldots,\sqrt{p_{6}Z_{6}}
\right].
\end{array}
\]

The vector of extended stress intensity factors (SIFs) is defined as
\begin{equation}
	\tilde{\mathbf{k}}^{(1)}
	=
	\left[
	K_{\mathrm{II}}^{\sigma},
	K_{\mathrm{I}}^{\sigma},
	K_{\mathrm{III}}^{\sigma},
	K_{\mathrm{II}}^{H},
	K_{\mathrm{I}}^{H},
	K_{\mathrm{III}}^{H}
	\right],
	\label{KushnirEq14}
\end{equation}
where the superscripts $\sigma$ and $H$ correspond to the phonon and
phason parts of the stress intensity factors, respectively, and the
subscripts I, II, and III denote the classical fracture modes:
opening (I), sliding (II), and tearing (III).

The extended stress intensity factors (SIFs) are defined based on the
asymptotic behavior of the extended stress field evaluated along the
crack extension line (i.e., ahead of the crack tip). Specifically,
\begin{equation}
	\tilde{\boldsymbol{\sigma}}_{2}(x'_1)
	=
	\frac{\tilde{\mathbf{k}}^{(1)}}{\sqrt{2\pi x'_1}}
	+ O(1),
	\label{KushnirEq15}
\end{equation}
where $x'_1$ is the local coordinate measured along the crack
extension line. The leading singular term is directly proportional
to the extended SIF vector $\tilde{\mathbf{k}}^{(1)}$, which
characterizes both the mechanical and thermal loading modes in
quasicrystalline solids.

The square-root singularity of the heat flux and extended stress
fields must be explicitly taken into account in the numerical solution
of the boundary integral equations. The corresponding discontinuity
functions defined on the crack surface (such as $\Delta\theta$ or
$\Delta\tilde{u}_{I}$) also exhibit square-root behavior near the
crack tips. Neglecting this feature would significantly reduce the
accuracy of the numerical solution.

Furthermore, the extended stress intensity factors serve as the key
fracture parameters in the mechanics of quasicrystals, in the same
way that classical SIFs characterize fracture behavior in conventional
materials. Their accurate evaluation is therefore essential for
predicting crack initiation and propagation and for performing
reliable fracture analysis of quasicrystalline solids.

\section{Boundary Element Method for Solution of the Obtained Equations}

The system of boundary integral equations
(\ref{KushnirEq7})--(\ref{KushnirEq10})
is solved using the boundary element method (BEM).
For this purpose, the boundary curves $\Gamma$ and $\Gamma_{C}$ are
discretized into $n$ and $n_{C}$ straight line segments, respectively,
called boundary elements $\Gamma_{q}$:
\[
\Gamma \cup \Gamma_{C}
\approx
\bigcup_{q=1}^{n+n_{C}} \Gamma_{q}.
\]

Each boundary element is defined by three nodal points: one located
at the center of the element and the other two positioned symmetrically
at one-third of the element length on either side of the center.
This leads to the use of discontinuous quadratic boundary elements
(also referred to as three-node discontinuous elements).

With this configuration, the collocation point—i.e., the point at
which the boundary integral equation is enforced—always lies on a
smooth portion of the piecewise-linear approximation of the boundary.
This ensures the correct enforcement of boundary conditions in the
discretized form of the integral equations.

The boundary functions $\theta$, $h_{n}$, $\Delta\theta$, $\Sigma h_{n}$,
$\tilde{u}_{I}$, $\tilde{t}_{I}$, and $\Delta\tilde{u}_{I}$ are
approximated within each element using their nodal values:
\begin{multline}
	\left[
	\theta,\, h_{n},\, \Delta\theta,\, \Sigma h_{n},\,
	\tilde{t}_{I},\, \tilde{u}_{I},\, \Delta\tilde{u}_{I}
	\right](\xi)
	\\
	\approx
	\sum_{p=1}^{3}
	\left[
	\theta^{q,p},\, h_{n}^{q,p},\, \Delta\theta^{q,p},\, \Sigma h_{n}^{q,p},\,
	\tilde{t}_{I}^{q,p},\, \tilde{u}_{I}^{q,p},\, \Delta\tilde{u}_{I}^{q,p}
	\right]
	\phi_{p}(\xi).
	\label{KushnirEq16}
\end{multline}
%
Here $\xi$ is the local coordinate on the element, varying within
the standard interval $-1 \leq \xi \leq 1$, and $\phi_{p}(\xi)$ are
the corresponding shape functions. The differential arc-length element
on $\Gamma_{q}$ is expressed as
\[
d\Gamma_{q} = J_{q}\, d\xi,
\]
where $J_{q}$ denotes the Jacobian (i.e., the modulus of the Jacobian
determinant) associated with the transformation from the global
coordinate system to the local element coordinate $\xi$.

Using this approach, the system of singular boundary integral equations
(\ref{KushnirEq7})--(\ref{KushnirEq10}) is reduced to a system of
linear algebraic equations with respect to the unknown nodal values
$\theta^{q,p}$, $h_{n}^{q,p}$, $\Delta\theta^{q,p}$, $\Sigma h_{n}^{q,p}$,
$\tilde{t}_{I}^{q,p}$, $\tilde{u}_{I}^{q,p}$, and $\Delta\tilde{u}_{I}^{q,p}$
of the boundary functions $\theta$, $h_{n}$, $\Delta\theta$, $\Sigma h_{n}$,
$\tilde{u}_{I}$, $\tilde{t}_{I}$, and $\Delta\tilde{u}_{I}$.

Thus, the resulting linear algebraic equations for plane heat conduction
in a cracked quasicrystal solid are obtained as follows.

\begin{itemize}
	
	\item For a collocation point $\mathbf{y}^{s,r}\in\Gamma$:
	
	\begin{multline}
		\frac{1}{2}\theta^{s,r}
		=
		\sum_{q=1}^{n}\sum_{p=1}^{3}
		\left[
		\Theta^{qpsr}h_{n}^{q,p}
		-
		H^{qpsr}\theta^{q,p}
		\right]
		\\
		+
		\sum_{q=n+1}^{n+n_{C}}\sum_{p=1}^{3}
		\left[
		\Theta^{qpsr}\Sigma h_{n}^{q,p}
		-
		H^{qpsr}\Delta\theta^{q,p}
		\right]
		\\
		-
		\iint_{S}
		\Theta^{*}(\mathbf{x},\mathbf{y}^{s,r})
		f_{h}(\mathbf{x})
		\,dS(\mathbf{x}).
		\label{KushnirEq17}
	\end{multline}
	
	\item For a collocation point $\mathbf{y}^{s,r}\in\Gamma_{C}$ on a crack
	surface $\Gamma_{C}$:
	
	\begin{multline}
		\frac{1}{2}\Delta h_{n}^{s,r}
		=
		\sum_{q=1}^{n}\sum_{p=1}^{3}
		\left[
		\Theta_{1}^{qpsr}h_{n}^{q,p}
		-
		H_{1}^{qpsr}\theta^{q,p}
		\right]
		\\
		+
		\sum_{q=n+1}^{n+n_{C}}\sum_{p=1}^{3}
		\left[
		\Theta_{1}^{qpsr}\Sigma h_{n}^{q,p}
		-
		H_{1}^{qpsr}\Delta\theta^{q,p}
		\right]
		\\
		-
		n_{i}^{+}(\mathbf{y}^{s,r})
		\iint_{S}
		\Theta_{i}^{**}(\mathbf{x},\mathbf{y}^{s,r})
		f_{h}(\mathbf{x})
		\,dS(\mathbf{x}).
		\label{KushnirEq18}
	\end{multline}
	
\end{itemize}
%
Here

\begin{equation}
	\begin{aligned}
		\Theta^{qpsr}
		&=
		\int_{-1}^{1}
		\Theta^{*}(\mathbf{x}(\xi),\mathbf{y}^{s,r})
		\phi_{p}(\xi)
		J_{q}(\xi)
		\,d\xi(\Gamma_{q}),
		\\ 
		H^{qpsr}
		&=
		\int_{-1}^{1}
		H^{*}(\mathbf{x}(\xi),\mathbf{y}^{s,r})
		\phi_{p}(\xi)
		J_{q}(\xi)
		\,d\xi(\Gamma_{q}),
		\\ 
		\Theta_{1}^{qpsr}
		&=
		n_{i}^{+}(\mathbf{y}^{s,r})
		\int_{-1}^{1}
		\Theta_{i}^{**}(\mathbf{x}(\xi),\mathbf{y}^{s,r})
		\phi_{p}(\xi)
		J_{q}(\xi)
		\,d\xi(\Gamma_{q}),
		\\ 
		H_{1}^{qpsr}
		&=
		n_{i}^{+}(\mathbf{y}^{s,r})
		\int_{-1}^{1}
		H_{i}^{**}(\mathbf{x}(\xi),\mathbf{y}^{s,r})
		\phi_{p}(\xi)
		J_{q}(\xi)
		\,d\xi(\Gamma_{q}).
	\end{aligned}
	\label{KushnirEq19}
\end{equation}

The same procedure yields the following linear algebraic equations
for plane thermoelasticity of quasicrystals.

\begin{itemize}
	
	\item For a collocation point $\mathbf{y}^{s,r}\in\Gamma$:
	
	\begin{multline}
		\frac{1}{2}\tilde{u}_{I}^{s,r}
		=
		\sum_{q=1}^{n}\sum_{p=1}^{3}
		\left[
		U_{IJ}^{qpsr}\tilde{t}_{J}^{q,p}
		-
		T_{IJ}^{qpsr}\tilde{u}_{J}^{q,p}
		+
		R_{I}^{qpsr}\theta^{q,p}
		+
		V_{I}^{qpsr}h_{n}^{q,p}
		\right]
		\\
		+
		\sum_{q=n+1}^{n+n_{C}}\sum_{p=1}^{3}
		\left[
		-
		T_{IJ}^{qpsr}\Delta\tilde{u}_{J}^{q,p}
		+
		R_{I}^{qpsr}\Delta\theta^{q,p}
		+
		V_{I}^{qpsr}\Sigma h_{n}^{q,p}
		\right]
		\\
		+
		\iint_{S}
		U_{IJ}(\mathbf{x},\mathbf{y}^{s,r})
		\tilde{f}_{J}(\mathbf{x})
		\,dS(\mathbf{x})
		-
		\iint_{S}
		V_{I}(\mathbf{x},\mathbf{y}^{s,r})
		f_{h}(\mathbf{x})
		\,dS(\mathbf{x}).
		\label{KushnirEq20}
	\end{multline}
	
	\item For a collocation point $\mathbf{y}^{s,r}\in\Gamma_{C}$:
	
	\begin{multline}
		\frac{1}{2}\Delta\tilde{t}_{I}^{s,r}
		=
		\sum_{q=1}^{n}\sum_{p=1}^{3}
		\left[
		D_{IJ}^{qpsr}\tilde{t}_{J}^{q,p}
		-
		S_{IJ}^{qpsr}\tilde{u}_{J}^{q,p}
		+
		Q_{I}^{qpsr}\theta^{q,p}
		+
		W_{I}^{qpsr}h_{n}^{q,p}
		\right]
		\\
		+
		\sum_{q=n+1}^{n+n_{C}}\sum_{p=1}^{3}
		\left[
		-
		S_{IJ}^{qpsr}\Delta\tilde{u}_{J}^{q,p}
		+
		Q_{I}^{qpsr}\Delta\theta^{q,p}
		+
		W_{I}^{qpsr}\Sigma h_{n}^{q,p}
		\right]
		\\
		+
		n_{i}^{+}(\mathbf{y}^{s,r})
		\bigg[
		\iint_{S}
		D_{IjK}(\mathbf{x},\mathbf{y}^{s,r})
		\tilde{f}_{K}(\mathbf{x})
		\,dS(\mathbf{x})
		-
		\iint_{S}
		W_{Ij}(\mathbf{x},\mathbf{y}^{s,r})
		f_{h}(\mathbf{x})
		\,dS(\mathbf{x})
		\bigg].
		\label{KushnirEq21}
	\end{multline}
	
\end{itemize}

Here

\begin{equation}
	\begin{aligned}
		U_{IJ}^{qpsr}
		&=
		\int_{-1}^{1}
		U_{IJ}(\mathbf{x}(\xi),\mathbf{y}^{s,r})
		\phi_{p}(\xi)
		J_{q}(\xi)
		\,d\xi(\Gamma_{q}),
		\\
		T_{IJ}^{qpsr}
		&=
		\int_{-1}^{1}
		T_{IJ}(\mathbf{x}(\xi),\mathbf{y}^{s,r})
		\phi_{p}(\xi)
		J_{q}(\xi)
		\,d\xi(\Gamma_{q}),
		\\ 
		V_{I}^{qpsr}
		&=
		\int_{-1}^{1}
		V_{I}(\mathbf{x}(\xi),\mathbf{y}^{s,r})
		\phi_{p}(\xi)
		J_{q}(\xi)
		\,d\xi(\Gamma_{q}),
		\\ 
		R_{I}^{qpsr}
		&=
		\int_{-1}^{1}
		R_{I}(\mathbf{x}(\xi),\mathbf{y}^{s,r})
		\phi_{p}(\xi)
		J_{q}(\xi)
		\,d\xi(\Gamma_{q}),
		\\ 
		D_{IJ}^{qpsr}
		&=
		n_{k}^{+}(\mathbf{y}^{s,r})
		\int_{-1}^{1}
		D_{IkJ}(\mathbf{x}(\xi),\mathbf{y}^{s,r})
		\phi_{p}(\xi)
		J_{q}(\xi)
		\,d\xi(\Gamma_{q}),
		\\ 
		S_{IJ}^{qpsr}
		&=
		n_{k}^{+}(\mathbf{y}^{s,r})
		\int_{-1}^{1}
		S_{IkJ}(\mathbf{x}(\xi),\mathbf{y}^{s,r})
		\phi_{p}(\xi)
		J_{q}(\xi)
		\,d\xi(\Gamma_{q}),
		\\ 
		Q_{I}^{qpsr}
		&=
		n_{k}^{+}(\mathbf{y}^{s,r})
		\int_{-1}^{1}
		Q_{Ik}(\mathbf{x}(\xi),\mathbf{y}^{s,r})
		\phi_{p}(\xi)
		J_{q}(\xi)
		\,d\xi(\Gamma_{q}),
		\\ 
		W_{I}^{qpsr}
		&=
		n_{k}^{+}(\mathbf{y}^{s,r})
		\int_{-1}^{1}
		W_{Ik}(\mathbf{x}(\xi),\mathbf{y}^{s,r})
		\phi_{p}(\xi)
		J_{q}(\xi)
		\,d\xi(\Gamma_{q}).
	\end{aligned}
	\label{KushnirEq22}
\end{equation}

This numerical discretization transforms the original boundary value
problem into a matrix form that can be efficiently solved using standard
linear algebra techniques.

For boundary elements that are not adjacent to the crack tips, standard
Lagrange polynomials are used as shape functions in the approximation
(\ref{KushnirEq16}). These functions are defined for a three-node
discontinuous boundary element with nodal positions
\[
\xi_{p}=\left[-\tfrac{2}{3},\,0,\,\tfrac{2}{3}\right].
\]
The corresponding shape functions are
\begin{equation}
	\phi_{1}=\xi\!\left(\frac{9}{8}\xi-\frac{3}{4}\right),\quad
	\phi_{2}=\left(1-\frac{3}{2}\xi\right)\!\left(1+\frac{3}{2}\xi\right),\quad
	\phi_{3}=\xi\!\left(\frac{9}{8}\xi+\frac{3}{4}\right).
	\label{KushnirEq23}
\end{equation}

To improve the accuracy of the method and to facilitate the extraction
of extended stress intensity factors, specialized three-node discontinuous
boundary elements are introduced for modeling the vicinity of the crack
tips.

According to the asymptotic relations (\ref{KushnirEq13}), the shape
functions for these special elements are designed to reproduce the
known crack-tip behavior.

\begin{itemize}
	
	\item For displacement or temperature discontinuities, the shape
	functions are chosen as
	\begin{equation}
		\phi_{p}^{\Delta}
		=
		\Phi_{p1}^{\Delta}\sqrt{\rho}
		+
		\Phi_{p2}^{\Delta}\rho
		+
		\Phi_{p3}^{\Delta}\rho^{3/2},
		\qquad p=1,2,3.
		\label{KushnirEq24}
	\end{equation}
	
	\item For heat flux jumps in the case of thermally conductive cracks,
	the following approximation is used:
	\begin{equation}
		\phi_{p}^{\Sigma}
		=
		\Phi_{p1}^{\Sigma}\rho^{-1/2}
		+
		\Phi_{p2}^{\Sigma}
		+
		\Phi_{p3}^{\Sigma}\sqrt{\rho},
		\qquad p=1,2,3.
		\label{KushnirEq25}
	\end{equation}
	
\end{itemize}
%
Here $\rho=1\pm\xi$, where the sign depends on whether the element
is located near the left or the right crack tip.

The matrices of constants $\Phi_{pj}^{\Delta}$ and $\Phi_{pj}^{\Sigma}$
are uniquely determined from the interpolation conditions
\begin{equation}
	\phi_{p}(\xi_{p})=1,
	\qquad
	\phi_{p}(\xi_{j\ne p})=0,
	\label{KushnirEq26}
\end{equation}
where $\xi_{p}$ are the local coordinates of the nodal points on the
boundary element.

These singular shape functions ensure that the approximation accurately
captures the square-root singularity characteristic of the physical
fields near the crack tips, which is essential for achieving high
accuracy in the numerical solution.

With the introduction of the singular shape functions
(\ref{KushnirEq24}) and (\ref{KushnirEq25}) into the boundary element
scheme (\ref{KushnirEq17})--(\ref{KushnirEq21}), it becomes necessary
to evaluate the following types of singular integrals:
\begin{equation}
	\begin{array}{rclrcl}
		
		I_{1} &=& \displaystyle\mathrm{RPV}\int_{-1}^{1} F(\ln\xi)\,
		\phi^{\Sigma}(\xi)\,d\xi,
		\quad &
		I_{2} &=& \displaystyle\mathrm{CPV}\int_{-1}^{1} F\!\left(\frac{1}{\xi}\right)
		\phi^{\Delta}(\xi)\,d\xi,
		\\[1em]
		I_{3} &=& \displaystyle\mathrm{CPV}\int_{-1}^{1} F\!\left(\frac{1}{\xi}\right)
		\phi^{\Sigma}(\xi)\,d\xi,
		\quad &
		I_{4} &=& \displaystyle\mathrm{HPV}\int_{-1}^{1} F\!\left(\frac{1}{\xi^{2}}\right)
		\phi^{\Delta}(\xi)\,d\xi .
	\end{array}
	\label{KushnirEq27}
\end{equation}

For integrals of type $I_{1}$, nonlinear transformations have proven
effective for smoothing the integrand in the vicinity of the
singularity. However, such transformations typically regularize the
integrand only at the singular point $\xi_{p}$ and do not eliminate
the end-point singularities that arise due to the specific form of
the shape functions (\ref{KushnirEq25}). To address this difficulty,
the integration interval $(-1,1)$ is first divided into two
subintervals at the singular point $\xi_{p}$. Each subinterval is then
normalized to the standard range $(0,1)$ so that the singular point
is mapped to the new origin $\xi=0$.

After normalization, each interval is mapped onto itself using the
nonlinear transformation
\begin{equation}
	\xi = 3s^{2}-2s^{3},
	\qquad
	d\xi = 6(s-s^{2})\,ds,
	\label{KushnirEq28}
\end{equation}
which smooths the integrand at both ends of the integration domain.
This transformation improves the accuracy and stability of the
numerical integration, particularly when singular shape functions are
used to model the behavior near crack tips.

To evaluate the other singular integrals $I_{2}$, $I_{3}$, and $I_{4}$,
special interpolation-type quadrature formulas based on Chebyshev nodes
are constructed. These quadratures approximate integrals interpreted
in the sense of the Cauchy principal value (CPV) or the Hadamard
finite part (HPV):
\begin{equation}
	\int_{-1}^{1}\frac{f(s)}{s}\,ds
	=
	\sum_{k=1}^{n}A_{k}^{\mathrm{CPV}}\,f(\sigma_{k}),
	\qquad
	\int_{-1}^{1}\frac{f(s)}{s^{2}}\,ds
	=
	\sum_{k=1}^{n}A_{k}^{\mathrm{HPV}}\,f(\sigma_{k}),
	\label{KushnirEq29}
\end{equation}
where the Chebyshev nodes are defined by
\[
\sigma_{k}
=
\cos\!\left(\frac{(2k-1)\pi}{2n}\right),
\]
and $n$ is a positive even integer.

These quadrature rules (\ref{KushnirEq29}) are exact when the integrand
$f(s)$ is a polynomial of degree less than $n$. The weights
$A_{k}^{\mathrm{CPV}}$ and $A_{k}^{\mathrm{HPV}}$ are determined by
substituting monomials of degree less than $n$ into the quadrature
formulas and equating the resulting expressions on both sides. This
procedure yields linear systems for the quadrature weights.

For the CPV integral we obtain the following system:
\begin{equation}
	\begin{aligned}
		\int_{-1}^{1}\frac{1}{s}\,ds
		&=
		\sum_{k=1}^{n}A_{k}^{\mathrm{CPV}}
		=
		0,
		\\
		\int_{-1}^{1}\frac{s^{j}}{s}\,ds
		&=
		\sum_{k=1}^{n}A_{k}^{\mathrm{CPV}}\sigma_{k}^{j}
		=
		\frac{1-(-1)^{j}}{j}
		\qquad
		(j=1,\ldots,n-1).
	\end{aligned}
	\label{KushnirEq30}
\end{equation}

For the HPV integral we obtain
\begin{equation}
	\begin{aligned}
		\int_{-1}^{1}\frac{s}{s^{2}}\,ds
		&=
		\sum_{k=1}^{n}A_{k}^{\mathrm{HPV}}\sigma_{k}
		=
		0,
		\\
		\int_{-1}^{1}\frac{s^{j}}{s^{2}}\,ds
		&=
		\sum_{k=1}^{n}A_{k}^{\mathrm{HPV}}\sigma_{k}^{j}
		=
		\frac{1+(-1)^{j}}{j-1}
		\qquad
		(j=0,2,3,\ldots,n-1).
	\end{aligned}
	\label{KushnirEq31}
\end{equation}

These systems can be solved numerically to determine the quadrature
weights with high accuracy. The resulting formulas allow accurate
and efficient evaluation of the singular integrals $I_{2}$, $I_{3}$,
and $I_{4}$.

The use of specially constructed shape functions
(\ref{KushnirEq24}) and (\ref{KushnirEq25}) enables significantly more
accurate evaluation of the phonon–phason (extended) stress intensity
factors, which appear as coefficients of the singular term
$\sqrt{\rho}$. These factors play a key role in characterizing the
coupled phonon–phason fields near singularities such as crack tips
or sharp inclusion edges.

For a terminal straight boundary element whose vertex corresponds to
$\xi=1$, the asymptotic relations (\ref{KushnirEq13}) together with
the displacement discontinuity shape functions
(\ref{KushnirEq24}) lead to the following computational formula for
the phonon–phason stress intensity factor:
\begin{equation}
	\tilde{\mathbf{k}}^{(1)}
	=
	\frac{\sqrt{\pi}}{4\sqrt{J}}
	\mathbf{L}
	\left(
	\Delta\tilde{\mathbf{u}}^{3-3'}\Phi_{11}^{\Delta u}
	+
	\Delta\tilde{\mathbf{u}}^{2-2'}\Phi_{21}^{\Delta u}
	+
	\Delta\tilde{\mathbf{u}}^{1-1'}\Phi_{31}^{\Delta u}
	\right),
	\label{KushnirEq32}
\end{equation}
where $\mathbf{L}$ is the extended Barnett–Lothe tensor for
quasicrystals, which accounts for the coupling between the phonon
and phason fields.

\section{Numerical Example}

Consider a square-shaped body with side length $2W$, composed of a
one-dimensional hexagonal quasicrystal. In this material, the atomic
structure is periodic in the $Ox_{1}x_{2}$ plane and quasiperiodically
arranged along the $Ox_{3}$ direction. The material properties are
listed in Table~\ref{KushnirTab1} \cite{Kushnir5}. The domain contains
a centrally located thermally insulated crack of length $2a$, as
shown in Figure~\ref{KushnirFig1}. The two sides of the square are
aligned parallel to the direction of quasiperiodic ordering, denoted
by the axis $Ox_{3}$, while the crack is oriented perpendicular to
this direction and lies along the axis $Ox_{1}$.

\begin{table}[h]
	\caption{Material constants for the 1D hexagonal quasicrystal \cite{Kushnir5}}
	\label{KushnirTab1}
	
	\centering
	\begin{tabular*}{\textwidth}{@{\extracolsep{\fill}}ll}
		\toprule
		Property & 1D hexagonal quasicrystal \\
		\midrule
		Phonon elastic moduli (GPa)
		& $c_{11}=150$,\;\; $c_{12}=55$,\;\; $c_{13}=45$,\;\; $c_{33}=90$,\;\; $c_{44}=50$ \\
		
		Phason elastic moduli (GPa)
		& $K_{1}=0.084$,\;\; $K_{2}=0.036$ \\
		
		Thermal moduli (MPa/K)
		& $\beta_{1}=1.798$,\;\; $\beta_{3}=1.383$ \\
		
		Phonon--phason coupling (GPa)
		& $R_{1}=-1.68$,\;\; $R_{2}=1.2$,\;\; $R_{3}=1.2$ \\
		
		Heat conduction (W/m/K)
		& $k_{11}=12.4$,\;\; $k_{33}=12.4$ \\
		\bottomrule
	\end{tabular*}
\end{table}

\begin{figure}[b]
	\sidecaption[b]
	\centering
	\includegraphics[width=.5\linewidth]{KushnirFig1}
	\caption{Sketch of the problem}
	\label{KushnirFig1}
\end{figure}

All four boundaries of the square are assumed to be mechanically free,
i.e., no external mechanical tractions are applied. On the two opposite
sides parallel to the quasiperiodic direction, constant values of
the normal components $h_{1}$ and $h_{3}$ of the heat flux density
vector are prescribed, simulating a steady-state thermal loading
across the domain.

Table~\ref{KushnirTab2} presents the values of the non-zero normalized
phonon–phason stress intensity factors at the right-hand crack tip
as functions of the crack-to-side length ratio $a/W$. These extended
stress intensity factors characterize the singular behavior of the
stress fields in the vicinity of the crack tip in a quasicrystalline
medium.

\begin{table}[h]
	\caption{Phonon--phason stress intensity factors for a thermally insulated crack
		in a square quasicrystal domain}
	\label{KushnirTab2}
	
	\centering
	\begin{tabular*}{\textwidth}{@{\extracolsep{\fill}}lccccccccc}
		\toprule
		${a/W}$ 
		& 0.1 & 0.2 & 0.3 & 0.4 & 0.5 & 0.6 & 0.7 & 0.8 & 0.9 \\
		\midrule
		
		\multicolumn{10}{c}{$h_{1}=0$, $h_{3}=h_{0}$} \\
		
		$-{K_{\mathrm{II}}^{\sigma}/K^{0}}$
		& 0.0946 & 0.0934 & 0.0919 & 0.0912 & 0.0923 & 0.0967 & 0.1069 & 0.1287 & 0.1841 \\
		
		\midrule
		
		\multicolumn{10}{c}{$h_{1}=h_{0}$, $h_{3}=0$} \\
		
		$10^{3}{K_{\mathrm{I}}^{\sigma}/K^{0}}$
		& 0.354 & 0.348 & 0.336 & 0.318 & 0.292 & 0.258 & 0.218 & 0.173 & 0.134 \\
		
		$10^{3}{K_{\mathrm{I}}^{H}/K^{0}}$
		& 0.038 & 0.040 & 0.043 & 0.048 & 0.055 & 0.064 & 0.076 & 0.093 & 0.122 \\
		
		\bottomrule
	\end{tabular*}
\end{table}

The normalization factor used for these extended SIFs is given by
\[
K^{0}=\frac{h_{0}\beta_{11}a\sqrt{\pi a}}{k_{11}}.
\]

In the boundary element model, the boundary of the square domain is
discretized into 80 uniformly distributed elements, while the crack
faces are represented by 20 elements, two of which are specially
designed to capture the singular thermal and mechanical behavior near
the crack tips.



When the heat flux is applied perpendicular to the crack surface
($h_{1}=0$, $h_{3}=h_{0}$), the thermally insulated crack disturbs
the thermal field in the body. Since heat cannot pass through the
crack, the heat flux is diverted around it, which leads to a
concentration of thermal stresses near the crack tips. For a short
crack ($a/W=0.1$), the obtained results are in good agreement with
the analytical solution for an infinite domain containing a single
crack \cite{Kushnir34}. For this case, the analytical solution gives
the normalized stress intensity factor
$K_{\mathrm{II}}^{\sigma}/K^{0}=-0.0949$.

As the ratio $a/W$ increases, the normalized phonon stress intensity
factor $K_{\mathrm{II}}^{\sigma}/K^{0}$ initially decreases and then
begins to increase after reaching a minimum value. This behavior is
explained by the fact that the normalization coefficient $K^{0}$
depends on the crack length. The actual (non-normalized) stress
intensity factors always increase with increasing crack length.
Only the phonon component $K_{\mathrm{II}}^{\sigma}$ is non-zero in
this case, while the first-mode (opening-mode) stress intensity
factors vanish.

In contrast, when the heat flux is applied along the crack
($h_{1}=h_{0}$, $h_{3}=0$), the crack does not disturb the thermal
field. Nevertheless, non-zero stress fields still arise due to
thermoelastic effects specific to quasicrystals
\cite{Kushnir43}. Such effects do not occur in ordinary crystalline
materials. In conventional crystals, the stress intensity factors
in this situation would be zero. However, in quasicrystals both
phonon and phason components become non-zero. The normalized phonon
component decreases with increasing $a/W$, whereas the phason
component increases. In addition, in this case only first-mode
stress intensity factors appear, in contrast to the previous case
where only second-mode components were observed.

\section{Conclusions}

An efficient boundary element method has been developed for the analysis
of quasicrystalline bodies containing cracks under thermomechanical
loading. A specialized numerical technique has been proposed for the
accurate computation of phonon--phason stress intensity factors. This
approach enables reliable extraction of generalized stress intensity
factors in the presence of singular thermal and mechanical fields near
crack tips.

The accuracy of the method has been validated through comparison with
available analytical solutions. In particular, excellent agreement was
obtained for small crack lengths, confirming the reliability of the
proposed formulation. The method exhibits good numerical stability and
convergence over a wide range of crack sizes. It therefore allows a
detailed investigation of the influence of crack length on the
phonon--phason stress intensity factors, which is especially important
in engineering applications where defect geometry and scale can
significantly affect the mechanical behavior of materials.

The proposed approach also extends the capabilities for analyzing
finite-sized quasicrystalline domains that more closely represent
practical systems. The method successfully captures the coupled thermal
and mechanical interactions in the vicinity of cracks. It is therefore
well suited for fracture mechanics analyses of quasicrystals, whose
distinctive thermoelastic behavior cannot be adequately described by
classical crystalline models. Furthermore, the approach provides
additional insight into thermomechanical phenomena specific to
quasicrystals, such as stress generation induced by a uniform
temperature gradient.


\bibliographystyle{unsrtnat}
\bibliography{Chapter24}
